\section{Related Work}\label{sec:related}

\subsection{Distributed statevector simulation on CPUs and GPUs}\label{subsec:related-dist}

Exact statevector simulation beyond the memory of one node distributes the $2^n$ amplitudes over many processes,
each holding a contiguous slice selected by the top qubits, and exchanges amplitudes when a gate acts on such a global
qubit. qHiPSTER, later the Intel Quantum Simulator, combined vectorisation, multithreading, cache blocking and
overlap of computation with communication and simulated circuits of up to 40 qubits on a supercomputer, limited by
memory and network bandwidth~\cite{smelyanskiy2016qhipster,guerreschi2020intel}. QuEST is a hybrid multithreaded,
distributed and GPU accelerated simulator that benchmarked random circuits of up to 38 qubits on up to 2048 nodes~\cite{jones2019quest}.
H\"aner and Steiger used generated and optimised kernels and a scheduling algorithm that reduces communication to
simulate a 45 qubit circuit on 8192 nodes~\cite{haner2017petabyte}; the J\"ulich simulator reached 48 qubits with an
adaptive encoding of the wave function and showed close to ideal weak scaling on several supercomputers~\cite{deraedt2019massively};
mpiQulacs is a distributed statevector simulator optimised for A64FX clusters~\cite{imamura2022mpiqulacs}. On GPUs,
HyQuas selects among partitioning methods per circuit pattern to improve data locality~\cite{zhang2021hyquas}, and the
cuQuantum SDK provides composable GPU primitives for quantum circuit simulation~\cite{bayraktar2023cuquantum};
Qulacs is a widely used single node CPU and GPU simulator~\cite{suzuki2021qulacs}. Simulations at 49 qubits and beyond
rely on other representations, such as tensor network contraction and low depth circuit
structure~\cite{markov2008simulating,pednault2017breaking,chen2018sixtyfour,boixo2018characterizing,villalonga2019flexible}.

Two techniques of these systems are central to the present work. First, a gate that is diagonal in the computational
basis never needs amplitudes of another process, even when it acts on a global qubit, and a gate controlled by a
global qubit acts only on the processes whose bit is one; second, global and local qubits can be swapped once so that
a sequence of gates on a global qubit runs locally (global to local remapping), and gates can be scheduled to reduce
the number of such exchanges~\cite{haner2017petabyte}. Executing diagonal gates on partition qubits without exchange is the
first technique; it is not new as an algorithmic idea, and we do not claim it as such. What this paper contributes is the
implementation of these techniques inside FPGA kernels: the slice enable lets every card's HBM engines execute
partition diagonal, partition controlled and partition permutation gates in their streaming passes without host
arithmetic or PCIe traffic, and remaps run as device copies inside a card and as half slice exchanges across cards.

\subsection{FPGA statevector simulators}\label{subsec:related-fpga}

FPGA emulation of quantum circuits was proposed by Khalid, Zilic and Radecka~\cite{khalid2004fpga}. Lee et al.\
presented a serial parallel FPGA emulation architecture with the QFT and Grover's search as case studies and a linear
reduction of resources compared with pipelined designs~\cite{lee2016fpga}; Pilch and D\l{}ugopolski proposed a
universal FPGA emulator with a standard programming interface~\cite{pilch2019fpga}; Mahmud, El-Araby and Caliga
studied scalable reconfigurable emulation at high precision and throughput~\cite{mahmud2019scaling}. More recent work
includes Q2Logic, a coarse grained reconfigurable overlay specialised for Schr\"odinger style
simulation~\cite{podobas2023q2logic}, PCQ, a parallel pipelined dataflow architecture with data decoupling that
partitions computation into non interacting subsets~\cite{pcq2024}, and a scalable accelerator architecture that
addresses parallelism and the routing of state components~\cite{belfore2024scalable}. These designs target one
device. Statevector simulation over several FPGAs has been demonstrated by the Trefoil system, which combines
multiple FPGAs with SATA disk arrays (2~PB of storage in eight storage subsystems) to hold state vectors of up to 47
qubits, at the cost of disk bandwidth (an H gate on 30 qubits takes about 10 minutes)~\cite{niwase2023trefoil}; its
simulator, Qu-Trefoil, builds on this storage based platform~\cite{wei2025qutrefoil}. The present work takes the
opposite design point: the whole state stays in the HBM of the cards and moves over PCIe only when a gate requires
amplitudes of another card, which limits the width to what the HBM of the cards holds but keeps every gate at HBM
bandwidth. The companion paper of this work~\cite{tiwari2026fpga30} reaches 30 qubits in complex64 on one U55C with HBM and
compares it with CPU and GPU simulators; the present paper is the multi card, complex128 continuation that the
companion paper names as future work. Further related designs of the author on the same platform, a native RTL
statevector engine in FP64~\cite{ali2026rtl}, real time error correction decoders~\cite{ali2026qec} and a matrix
product state engine~\cite{ali2026mps}, are available as preprints.
